\section{Artin and Shephard applications}\label{sec:app-artin}
The common cyclic-marking theorem proves all intersection, coset Helly, residue and panel assertions after an exact semidirect-product identification. Its two presentation theorems supply inverse maps, rather than only representations of a quotient. We retain the classical standard-parabolic inputs for Artin and Coxeter groups, with their exact references stated in the common paper. General Shephard groups do not receive those stronger properties merely from their presentations. The following separate argument supplies exactly the generator and diagram-action statement needed for source Lemma~5.6.

\subsection{Separating the named Shephard generators}\label{S-sec:reflection}

A properly labelled Shephard graph has a finite vertex set $S$, vertex orders
$p_i\in\{2,3,\ldots\}\cup\{\infty\}$, and pair labels
$m_{ij}\in\{2,3,\ldots\}\cup\{\infty\}$. An odd finite pair label
requires $p_i=p_j$. Its group has presentation
\[
 \langle s_i\ (i\in S)\mid s_i^{p_i}=1\ (p_i<\infty),\quad
 \W_{m_{ij}}(s_i,s_j)=\W_{m_{ij}}(s_j,s_i)\ (m_{ij}<\infty)\rangle,
\]
where $\W_m$ denotes an alternating word of length $m$. A pair of label two
commutes, while an infinite label imposes no pair relation. All vertex
orders infinite give the Artin presentation; all vertex orders two give
the Coxeter presentation. Graph automorphisms preserve both sets of labels.
We now close the injectivity point in source Lemma~5.6 without assuming
an embedding theorem for general Shephard groups. This argument constructs
a representation separating the named generators; it does not assert that
the representation is faithful on the entire group.

\begin{proposition}\label{S-prop:shephard}
In the presentation of Definition~5.1, every named generator has its
prescribed order and distinct names give distinct elements. Consequently
the action of the labelled graph automorphism group on the presented
Shephard group is faithful.
\end{proposition}
\begin{proof}
Let $V$ have basis $(e_i)_{i\in S}$ over $\mathbb C$. For vertex order
$p_i<\infty$ choose $z_i=\exp(2\pi\mathrm i/p_i)$, and for infinite
order choose $z_i=2$. Thus $z_i\ne1$ in every case, and equal vertex
orders give equal $z_i$. Fix square roots $u_i^2=z_i$.
Define linear maps
\[
 R_i(e_i)=z_i e_i,\qquad R_i(e_j)=e_j+c_{ij}e_i\quad(j\ne i).
\]
Choose the two coefficients independently for each unordered pair,
after ordering its endpoints as $i<j$. For a commuting pair take
$c_{ij}=c_{ji}=0$; for label infinity do the same, since no braid
needs to be imposed. For finite labels greater than two take
\[
 c_{ij}=1,\qquad c_{ji}=
 \begin{cases}
 -4z_i\cos^2(\pi/m),&m\text{ odd},\\
 2u_i u_j\cos(\pi/k)-z_i-z_j,&m=2k.
 \end{cases}
\]
Proper labelling ensures $z_i=z_j$ in the odd case.
The maps are invertible, with determinant $z_i$. The geometric-series
formula for their exceptional row shows $R_i^{p_i}=1$ for finite $p_i$;
its eigenvalue $z_i$ gives exact order $p_i$. If $p_i=\infty$, the
eigenvalue two gives infinite order.

We verify the braid relations on the whole vector space. On
$V_{ij}=\langle e_i,e_j\rangle$, put
\[
 A=\begin{pmatrix}z&c\\0&1\end{pmatrix},\qquad
 B=\begin{pmatrix}1&0\\d&w\end{pmatrix},\qquad P=AB.
\]
For $m=2k$, the trace and determinant of $P$ are
$2uv\cos(\pi/k)$ and $(uv)^2$, where $z=u^2,w=v^2$.
Its distinct eigenvalues are $uv\exp(\pm\pi\mathrm i/k)$.
Thus $P^k=-(uv)^kI=(BA)^k$, the required braid.

For $m=2r+1$, put $z=w$, $q=-\cos(2\pi/m)$ and use
$cd=-4z\cos^2(\pi/m)$. Then $\operatorname{tr}P=2zq$ and
$\det P=z^2$. Let $U_{-1}=0$, $U_0=1$ and
$U_{j+1}(q)=2qU_j(q)-U_{j-1}(q)$. Cayley--Hamilton gives
\[
 P^r=z^{r-1}U_{r-1}(q)P-z^rU_{r-2}(q)I.
\]
Direct multiplication also gives
$PA-BP=(z+cd)(A-B)$. Since $1+cd/z=2q-1$, subtraction yields
\[
 P^rA-BP^r=z^r\bigl(U_r(q)-U_{r-1}(q)\bigr)(A-B)=0.
\]
For the last equality write $q=\cos\theta$ with
$\theta=\pi-2\pi/m$. The recurrence has solution
$U_j(q)=\sin((j+1)\theta)/\sin\theta$; hence the difference is
zero because $\cos(m\theta/2)=0$. The two sides before subtraction
are precisely the alternating words of length $2r+1$.

To pass from $V_{ij}$ to $V$, set
$D=(z-1)(w-1)-cd$. When $D\ne0$, the common fixed space of
$R_i,R_j$ is complementary to $V_{ij}$: solve its two fixed-vector
linear equations by adjusting the $e_i,e_j$ coordinates. Both maps
are the identity on that complement, so the verified pair relation
holds on $V$. The exceptional values $D=0$ cause no gap. In the odd
case vary $z$ while keeping all other row coefficients fixed and
$cd=-4z\cos^2(\pi/m)$; $D$ is a nonzero polynomial in $z$.
Every entry of the braid difference is a polynomial vanishing at all
but finitely many complex $z$, and thus vanishes also at the exceptional
values. In the even case vary $u$ with the nonzero $v$ fixed;
$D=u^2v^2+1-2uv\cos(\pi/k)$ is again a nonzero polynomial.
The same polynomial argument applies, including the finitely many
excluded zero values. For a commuting pair the off-diagonal pair
coefficients are zero and the two exceptional rows commute directly.
All defining relations therefore hold on $V$.

The presentation now gives a homomorphism to $\operatorname{GL}(V)$.
If $i\ne j$, the $i$th diagonal entry of $R_i$ is $z_i\ne1$,
whereas that of $R_j$ is one. Thus these matrices, and consequently
the named group elements, are distinct. A labelled graph permutation
and its inverse preserve all defining relators, so they induce inverse
group automorphisms. A nonidentity graph permutation moves a named
generator to a different one, proving faithfulness.
\end{proof}
The constructed representation is used only to separate names and
certify their orders. No parabolic embedding, intersection or coset
Helly property for arbitrary Shephard groups follows from it.


\subsection{Free products and quotient actions}
If $A,B$ are Artin or Coxeter groups with disjoint standard markings,
$A*B$ has the union presentation, equivalently all cross labels are
infinite. Its generated subgroup on $R$ is
$A(R\cap S_A)*B(R\cap S_B)$ by reduced normal forms. Two such
subgroups intersect by intersecting the two factor markings: the
unique alternating reduced word belongs to both exactly when each
syllable does. The coset Helly theorem for free products proved in
the common paper
then applies to the two classical input geometries. Its hypotheses
allow arbitrary Borel subgroups; here both Borels are trivial.
The resulting standard marking therefore satisfies
the common parabolic-marking lemma. This recovers source Theorem~5.4 and all
three clauses of Corollary~5.5, for arbitrary choices of the two
Artin/Coxeter factors. Each source panel size is preserved.


\subsection{The involutory and cyclic families}
Apply the common involutory presentation theorem to two leaves with equal labels to every remaining vertex. This gives source Theorems~5.14--5.15 for every allowed finite or infinite label, with all their parabolic product and flag-complex conclusions. Its special $D_n$ case fixes the diagram convention explicitly:

For type $D_n$ ($n\geq4$), take the chain $a_1,a_2,\ldots,a_{n-1}$
with consecutive label three, attach $a_n$ by label three to $a_2$,
and let all other distinct pairs commute. Interchange $a_1,a_n$.
The new marking $t,a_1,a_2,\ldots,a_{n-1}$ has vertex orders
$2,\infty,\ldots,\infty$, consecutive labels $4,3,\ldots,3$ and
all nonconsecutive commutations. It is exactly $B_n(2,\infty)$.
This proves source Propositions~5.12--5.13; the arbitrary labels in the common involutory theorem
prove Theorems~5.14--5.15, including every stated intersection and
flag-complex assertion. The $t$ panels have size two and all remaining
panels are infinite.

There are two literal presentation qualifications. On source page~46
the graph attaches $a_n$ to $a_2$, whereas the printed relator list
uses $a_3$ and does not consistently bound the chain index. The proof
on page~47 uses the graph's $a_2$ convention, as do we. Also the
$B_n(p,q)$ display on page~41 omits the nonconsecutive commutations;
they are required by Definition~5.1 and by the displayed graph. They
are explicitly retained here. These are corrected graph presentations,
not verification of the malformed relator lists as printed.


For cyclic amalgamation take $k$ copies of the Artin group on $Y\sqcup\{x\}$ over the full standard parabolic on $Y$. The universal property identifies the amalgam $M$ with the Artin presentation having $k$ leaves and no relations between distinct leaves. The common cyclic presentation theorem gives exactly $M\rtimes C_k$, with vertex $t$ of order $k$, no relation on $x_0,t$, and all mixed braids inherited from one copy. Its complete inverse maps and cyclic-marking proof supply every geometric conclusion, with panel sizes $k$ and infinity.

These are source Theorems~5.16--5.17. The source's literal statement
writes $A_0\rtimes C_k$ where its construction and proof require
$M\rtimes C_k$. There need not even be the indicated cyclic action
on the single copy $A_0$. Its page~51 preliminary braid display also
has identical words on both sides; the right side must start with
$y$. When $k=2$, the order-two vertex is represented by no loop under
Definition~5.1, rather than by a loop labelled two. None of these
conventions affects the corrected presentation just proved.


\subsection{The tetrahedral Artin example}
In Example 5.7 let $A$ be the Artin group on $a_1,a_2,a_3,a_4$
with all pair labels three. The permutations $p=(23)$ and $q=(34)$ generate a group $B$ acting on the names. At representative 2 the orbits of
$1,\langle q\rangle,\langle p\rangle,B$ are
$\{2\},\{2\},\{2,3\},\{2,3,4\}$. These give the exact
orbit-intersection condition; the fixed singleton orbit $\{1\}$
also satisfies it. The complete orbit-transfer theorem applies.
The maximal parabolics are
\[
 A(\{2,3,4\})\rtimes B,\quad A(\{1\})\rtimes B,\quad
 A(\{1,2\})\rtimes\langle q\rangle,\quad
 A(\{1,2,3\})\rtimes\langle p\rangle.
\]
They are the deleted-generator subgroups of $(a_1,a_2,p,q)$.
This proves the source's firm, residually connected, flag-transitive
rank-four geometry; its panels have sizes $\infty,\infty,2,2$.

There is also a four-generator presentation, with $a=a_1,x=a_2$:
\[
 \begin{gathered}
 p^2=q^2=(pq)^3=1,\qquad [a,p]=[a,q]=[x,q]=1,\\
 \W_3(a,x)=\W_3(x,a),\qquad\W_6(x,p)=\W_6(p,x).
 \end{gathered}
\]
Set $x_2=x$, $x_3=pxp$, $x_4=qpxpq$. The relation $pqp=qpq$
and $[x,q]=1$ show that $p$ swaps $x_2,x_3$ and fixes $x_4$,
while $q$ swaps $x_3,x_4$ and fixes $x_2$.
The doubled-braid identity converts the length-six relation into
the length-three braid on $x_2,x_3$. Conjugation supplies the other
pair braids and the braids with $a$. These formulas give mutually
inverse maps with the original six-generator semidirect presentation.
